Remaining useful life (RUL) estimates are consumed by maintenance decisions that weigh premature replacement against unplanned failure, yet the literature still reports RUL almost exclusively as a point estimate. Interval estimates exist, but they split into two unsatisfactory families: distributional predictors that offer no finite-sample coverage guarantee, and conformal post-processing applied to a frozen predictor, whose interval width is decoupled from the asset's actual state. This paper shows that the second family has a hard ceiling on run-to-failure data, and proposes a method that removes it: the predictor is trained end-to-end under a dual-quantile loss, so that conditional width is learned rather than fitted post hoc, and the resulting interval is calibrated with conformal quantile regression. On four C-MAPSS subsets with engine-disjoint folds (80 training cells, four seeds), the method wins 75 of 80 paired comparisons at matched coverage of 0.90 (median 10.8\\% lower Winkler score; bootstrap 95\\% CI for the mean gain [7.4\\%, 10.6\\%]), and 19 of 20 at the nominal level (median 8.7\%). The gain holds on every subset, from 14.1\\% on FD002 to 6.5\\% on FD003. A cross-domain study on the FEMTO/PRONOSTIA bearing dataset reproduces the ranking on five of six bearings (mean 12.2\%). We further show that reporting mean coverage is insufficient: the worst engine falls below 0.80 coverage in 80--100\% of folds.
